\documentclass[10pt,a4paper]{amsart}
\usepackage[margin=19mm]{geometry}
\usepackage{amsmath,amssymb,mathtools}
\usepackage[scr=rsfs,cal=euler]{mathalfa}
\usepackage{xfrac}
\usepackage[unicode,hidelinks,bookmarks=false]{hyperref}
\newcommand{\ie}{i.e.\ }
\newcommand{\cf}{cf.\ }
\newcommand{\resp}{respectively\ }
\newcommand{\Subsection}{\subsection}
\providecommand{\nobreakdash}{\nobreak\hbox{-}}
\setlength{\emergencystretch}{2em}
\multlinegap0pt
\usepackage{bm}
\usepackage{bbm}
\usepackage{dsfont}
\usepackage{xspace}
\usepackage{cancel}
\usepackage{mathtools}
\usepackage{amsthm}
\makeatletter
\@ifundefined{theorem}{\newtheorem{theorem}{Theorem}}{}
\@ifundefined{lemma}{\newtheorem{lemma}[theorem]{Lemma}}{}
\makeatother
\makeatletter
\DeclareMathOperator{\Sym}{\bbR^{2\times 2}_{\rm sym}}
\newcommand{\bbR}{\mathbb{R}}
\newcommand{\calD}{\mathcal{D}}
\newcommand{\calL}{\mathcal{L}}
\newcommand{\calM}{\mathcal{M}}
\newcommand{\calX}{\mathcal{X}}
\expandafter\ifx\csname claim\endcsname\relax
\newenvironment{claim}[2]{\par\noindent\underline{Claim #2:}\space#1}{}{}
\fi
\newcommand\e{\varepsilon}
\newcommand\embed{\hookrightarrow} % embeddding
\newcommand{\mres}{\mathbin{\vrule height 1.6ex depth 0pt width
		0.13ex\vrule height 0.13ex depth 0pt width 1.3ex}}
\newcommand\ringring[1]{%
	{% make an Ord atom
		\mathop{\kern0pt #1}\limits^{% set a box over the variable
			\vbox to-1.85ex{
				\kern-2ex % lower the ring accents
				\hbox to 0pt{\hss\normalfont\kern.1em \r{}\kern-.45em \r{}\hss}%
				\vss % fill
			}% end of \vbox
		}% end of the superscript
	}% end of \mathop
}
\newcommand\weak{\rightharpoonup} % weak
\newcommand\weakstar{\stackrel{\ast }{\rightharpoonup}}  % weakstar
\makeatother
\title[Derivations of two one-dimensional models for transversely c]{Derivations of two one-dimensional models for transversely curved shallow shells: one leads to relaxation}
\author{Paroni Roberto, Picchi Scardaoni Marco}
\date{Source: arXiv:2507.23545v1 (CC BY 4.0)}
\begin{document}
\maketitle

\noindent\emph{Source and licence.} Derivations of two one-dimensional models for transversely curved shallow shells: one leads to relaxation. Version arXiv:2507.23545v1 (CC BY 4.0), distributed under the Creative Commons Attribution 4.0 International licence, as recorded in the arXiv licence metadata for this exact version and shown on the abstract page https://arxiv.org/abs/2507.23545v1. Full rights notice: licenses/arxiv-2507.23545-CC-BY-4.0.txt.\par\smallskip

\noindent\textit{Editorial scope.} Counted unit: the Lemma whose source label is \texttt{lemmawsq} in the article (source0.tex lines 825--833, source label \texttt{lemmawsq}), delivered together with the article's own complete original proof of it (source0.tex lines 834--903, one \texttt{proof} environment of 70 lines). No mathematics is written, repaired or re-derived: statement and proof are the article's own text line for line. Required context retained uncounted: cmp00 (lines 483--485); compactness4 (lines 496--528); cmp1 (lines 501--514); cmp2 (lines 504--508); cmp3 (lines 510--513); Sineq (lines 620--622); eqdet (lines 624--626); lemmaunif (lines 784--787); lemmamu2 (lines 812--815); cmp4 (lines 814--814); lemmader2 (lines 1592--1596). Every definition, display and label of those lines is retained unchanged. Omitted from this package and not redistributed: 1--482, 486--495, 509--509, 514--619, 623--623, 627--783, 788--811, 815--824, 904--1591, 1597--1980. Nothing inside a retained excerpt is omitted or altered: every retained body is delivered exactly as the article's lines are written, so no omission receipt of any kind accompanies this package. Citations: the retained excerpts carry no citation command at all, so no bibliography entry of the article is transcribed and no reference is redistributed. Reference adaptation: none was needed and none was made; every reference inside the delivered unit resolves inside it, so no unresolved reference or citation is produced. Printed slips kept literal: no printed word, symbol, hypothesis or number of the counted statement or of its proof was altered. Layout and class: the article's own class is \texttt{amsart} with packages including geometry, amssymb, mathtools, hyperref, enumerate, natbib, doi, tikz, tikz-3dplot, showkeys, pgfplots, lineno; the wrapper is the lane's pinned \texttt{amsart} editorial wrapper at 10pt on A4, which restates the theorem-like environments the retained text needs and adds the article's own macro definitions that the retained text needs (\texttt{\textbackslash Sym}, \texttt{\textbackslash bbR}, \texttt{\textbackslash calD}, \texttt{\textbackslash calL}, \texttt{\textbackslash calM}, \texttt{\textbackslash calX}, \texttt{\textbackslash e}, \texttt{\textbackslash embed}, \texttt{\textbackslash mres}, \texttt{\textbackslash ringring}, \texttt{\textbackslash weak}, \texttt{\textbackslash weakstar}), with extra wrapper packages (bm, bbm, dsfont, xspace, cancel, mathtools). The delivered numbering is the unit's own. Assets: no retained span contains a figure environment, a graphics-inclusion command, a PDF-page inclusion, a table or a TikZ picture; no image file is copied into this package, the package declares no asset dependency and no image file is redistributed with it. Licence: version arXiv:2507.23545v1 of this article is distributed under the Creative Commons Attribution 4.0 International licence, as recorded in the arXiv licence metadata for this exact version and shown on the abstract page https://arxiv.org/abs/2507.23545v1. A full rights notice is delivered at licenses/arxiv-2507.23545-CC-BY-4.0.txt. Characters: every retained excerpt is pure ASCII, so the delivered engine, XeLaTeX with its default Latin Modern text font, reports no missing character.
	\begin{equation}\label{cmp00}
		\sup_{\e} F_{\e}^{\beta}(u_\e, w_{\e}) <\infty.
	\end{equation}

	\begin{lemma}\label{compactness4}
		Let $0< \beta \le 2$.
		Let $(u_\e, w_{\e})\subset \calX$ be a sequence that satisfies \eqref{cmp00}.
		Then there exist $r \in H^1(I)$, $\vartheta \in H^1(I)$,   $\gamma_{11}, \gamma_{22}\in L^2(\Omega)$
		such that, up to a subsequence,
		\begin{align}
			&\partial_{2} \frac{w_\e}{\e} - \mathring{w} ' \weak \vartheta &&in \ H^1(\Omega), \label{cmp1}\\
			&\nabla^2_\e w_\e - \frac{\mathring{w} ''}{\e}\mathsf{e}_2\otimes \mathsf{e}_2 \weak
			\begin{pmatrix}
				\gamma_{11} &  \vartheta'
				\\
				\vartheta' & \gamma_{22}
			\end{pmatrix}
			&& in \ L^{2}(\Omega, \Sym), \label{cmp2}\\
			&\frac{w_\e}{\e^{\beta/2}} \weak w=\begin{cases}
				r & \text{if }\beta <2\\
				r+x_2\vartheta+\mathring{w} & \text{if }\beta =2
			\end{cases}  &&in \ H^1(\Omega).\label{cmp3}
		\end{align}
		Also, there exist $S \in L^2(\Omega,\mathbb{R}^{2\times 2}_{\rm sym})$, $(\xi_1, \xi_2) \in BV(I)\times BH(I)$,  such that, up to a subsequence,
		\begin{align}
			S_\e &\weak S&& in \ L^{2}(\Omega, \Sym), \label{cmp5}\\
			\frac{u_{\e}}{\e^\beta} &\weakstar u &&in \ BD(\Omega, \bbR^2),\label{cmp6}
		\end{align}
		with $u$ given by
		\begin{equation}\label{cmp7}
			u = \begin{cases}
				(\xi_1 - x_2 \xi_2', \xi_2) & \text{ if } \beta < 2\\
				(\xi_1 -x_2(\xi_2'+r'\vartheta + c_1\vartheta') - \mathring{w} r'   -\ringring{w}_{\langle 0 \rangle}\vartheta',\\
				\hspace{5cm}		\xi_2 - \frac12 x_2 \vartheta^2 - \mathring{w}\vartheta) & \text{ if } \beta = 2.
			\end{cases}
		\end{equation}
	\end{lemma}

		\begin{multline}\label{Sineq}
			\Big|\int_\Omega\varphi \det \nabla^2\frac{w_\e}{\e^{\beta/2}} \,dx\Big|\le \\ C\big(\|\partial_{22}\varphi\|_{L^2(\Omega)}+\e\|\partial_{12}\varphi\|_{L^2(\Omega)}+\e^2\|\partial_{11}\varphi\|_{L^2(\Omega)}\big) \qquad  \forall \ \varphi\in H^2_0(\Omega).
		\end{multline}

		\begin{equation}\label{eqdet}
			\det \nabla^2\frac{w_\e}{\e^{\beta/2}}=\e^{2-\beta}d_\e+\e^{1-\beta/2}\partial_{11}\frac{w_\e}{\e^{\beta/2}}\mathring{w}''
		\end{equation}

	\begin{lemma}\label{lemmaunif} Let $\beta =2$.
		Let $(u_\e, w_{\e})\subset \calX$  as in Lemma \ref{compactness4}. Then,
		$\frac{w_\e}{\e} \to w$ in $L^\infty(\Omega)$.
	\end{lemma}

	\begin{lemma}\label{lemmamu2}
		Let $\beta=2$. Let $(u_\e, w_\e)\subset \calX$ as in Lemma \ref{compactness4}. Then, there exists $\mu \in \calM^+(\Omega)$ such that
		\begin{equation}\label{cmp4}(\partial_1 \frac{w_\e}{\e^{\beta/2}})^2 \calL^2 \weakstar (\partial_1 w)^2\calL^2 + 2\mu \qquad  in \ \calM(\Omega).\end{equation}
	\end{lemma}

	\begin{lemma}\label{lemmader2}
		Let $f \in H^1(I)$. There exists a sequence $(f_n) \subset C^\infty(\overline{I})$ such that $f_n \to f$ in $H^1(I)$ as $n \uparrow \infty$. Moreover, if $(\e_n)\subset \bbR$ is such that $\e_n \downarrow 0$ as $n \uparrow \infty$,  $(f_n)$ can be chosen such that
		$\e_n f_n''\to 0$ in $L^2(I)$.
		%and $\e_n f_n' \to 0$ in $L^\infty(I)$ simultaneously.
	\end{lemma}

	\begin{lemma}\label{lemmawsq}
		Let $\beta =2$.
		Let $\mu$ as in Lemma \ref{lemmamu2}. Then, the support of $\mu$ is contained in $I\times Y$.
		In particular,
		for  every $y\in Y$ there is $\lambda_y \in \calM^+(I\times \{y\})$ such that
		$$
		\mu = \sum_{y\in Y} \lambda_y\otimes \delta_y.
		$$
	\end{lemma}

	\begin{proof}
		Let $w=r+x_2\vartheta+\mathring{w}$ be as in \eqref{cmp3}.
		As $r,\vartheta \in H^1(I)$, by Lemma \ref{lemmader2} there are  sequences $(r_\e), (\vartheta_\e) \subset C^\infty(\overline{I})$ such that $(r_\e, \vartheta_\e) \to (r, \vartheta)$ in $H^1(I)\times H^1(I)$  and $(\e r_\e'',\e \vartheta_\e'')\to (0,0)$ in $L^2(I)\times L^2(I)$.  The convergence $(r_\e, \vartheta_\e) \to (r, \vartheta)$ is actually uniform by the compact embedding $H^1(I) \embed C^0(\overline{I})$. Thus, the sequence $\tilde w_\e:=r_\e + x_2 \vartheta_\e+\mathring{w}$ converges uniformly to $w$ in $\Omega$.
		Combining this with Lemma \ref{lemmaunif}, we have that
		$z_\e\coloneqq   \frac{w_\e}{\e} - \tilde w_\e$ converges uniformly to $0$ in $\Omega$.
		By using \eqref{cmp3} and \eqref{cmp1} we deduce that
		\begin{equation}\label{ze1}
			\begin{aligned}
				\e\|\partial_1 z_\e\|_{L^2(\Omega)}&=\e\|\partial_1 \frac{w_\e}\e-r_\e'-x_2\vartheta_\e'\|_{L^2(\Omega)}\to 0,\\
				\|\partial_2 z_\e\|_{L^2(\Omega)}&=\|\partial_2 \frac{w_\e}\e-\vartheta_\e-\mathring{w}'\|_{L^2(\Omega)}\to 0
			\end{aligned}
		\end{equation}
		as $\e$ goes to zero. Also, by using \eqref{cmp2} we find that
		\begin{equation}\label{ze2}
			\begin{aligned}
				\e^2\|\partial_{11} z_\e\|_{L^2(\Omega)}&=\e\|\partial_{11} w_\e-\e r_\e'-x_2\e\vartheta_\e'\|_{L^2(\Omega)}\to 0,\\
				\e\|\partial_{12} z_\e\|_{L^2(\Omega)}&=\e\|\partial_{12} \frac{w_\e}\e-\vartheta_\e'\|_{L^2(\Omega)}\to 0,\\
				\|\partial_{22} z_\e\|_{L^2(\Omega)}&=\|\partial_{22} \frac{w_\e}\e-\mathring{w}''\|_{L^2(\Omega)}\to 0.
			\end{aligned}
		\end{equation}
		Let $\psi \in C_c^\infty(\Omega)$. By taking
		$\varphi_\e = \psi z_\e$ in place of $\varphi$ in \eqref{Sineq} we have
		\begin{equation}\label{Sineq2}
			\Big|\int_\Omega\varphi_\e \det \nabla^2\frac{w_\e}{\e} \,dx\Big|\le C\big(\|\partial_{22}\varphi _\e\|_{L^2(\Omega)}+\e\|\partial_{12}\varphi_\e\|_{L^2(\Omega)}+\e^2\|\partial_{11}\varphi_\e\|_{L^2(\Omega)}\big).
		\end{equation}
		From \eqref{ze1} and \eqref{ze2} it immediately follows that the left side of \eqref{Sineq2} goes to zero as $\e\to 0$.

		We now study the right side of \eqref{Sineq2}. By \eqref{eqdet}
		we have
		\begin{align*}
			0=\lim_{\e\to 0}\int_\Omega\varphi_\e \det \nabla^2\frac{w_\e}{\e}\,dx&=\lim_{\e\to 0}\int_\Omega\psi z_\e \big(d_\e+  \mathring{w}''\partial_{11}\frac{w_\e}{\e}\big)\,dx\\&=\lim_{\e\to 0}\int_\Omega\psi z_\e   \mathring{w}''\partial_{11}\frac{w_\e}{\e}\,dx,
		\end{align*}
		where the last equality follows since $d_\e$ is bounded in $L^1(\Omega)$ and $z_\e$ converges uniformly to zero.
		But by \eqref{cmp3} we deduce that
		\begin{multline}
			0=-\lim_{\e\to 0}\int_\Omega\psi z_\e   \mathring{w}''\partial_{11}\frac{w_\e}{\e}\,dx=\lim_{\e\to 0}\int_\Omega\big(z_\e\partial_{1}\psi    +\psi\partial_1z_\e\big)\mathring{w}''\partial_{1}\frac{w_\e}{\e}\,dx \\
			=
			\lim_{\e\to 0}\int_\Omega\psi\mathring{w}''\partial_1z_\e\partial_{1}\frac{w_\e}{\e}\,dx,
		\end{multline}
		that is equivalent to
		$$
		\lim_{\e\to 0}\int_\Omega\psi\mathring{w}''(\partial_1\frac{w_\e}{\e})^2\,dx
		=
		\lim_{\e\to 0}\int_\Omega\psi\mathring{w}''\partial_1\tilde w_\e\partial_{1}\frac{w_\e}{\e}\,dx=\int_\Omega\psi\mathring{w}''(\partial_{1}w)^2\,dx,
		$$
		as $\tilde w_\e\to w$ in $H^1(\Omega)$ and $w_\e/\e\weak w$ in $H^1(\Omega)$.
		So $\mathring{w}''(\partial_1\frac{w_\e}{\e})^2 \to \mathring{w}''(\partial_1 w)^2$ in $\calD'(\Omega)$.
		Since the sequence $(\mathring{w}''(\partial_1\frac{w_\e}{\e})^2)$ is bounded in $L^1(\Omega)$, we have that
		$$\mathring{w}''(\partial_1 \frac{w_\e}{\e})^2 \calL^2 \weakstar \mathring{w}''(\partial_1 w)^2\calL^2  \quad \mbox{in } \calM(\Omega).
		$$
		On the other hand, by \eqref{cmp4}, we also have that
		$$\mathring{w}''(\partial_1 \frac{w_\e}{\e})^2 \calL^2 \weakstar \mathring{w}''(\partial_1 w)^2\calL^2 + 2\mathring{w}''\mu \quad \mbox{in } \calM(\Omega).
		$$
		Hence, by uniqueness of the weak limits, we deduce that $\mathring{w}''\mu=0$.
		This implies that also the total variation of the measure $\mathring{w}''\mu$ is null, that is
		$|\mathring{w}''\mu|=|\mathring{w}''|\mu=0$. We now show that $\mu$ concentrates on $I\times Y$.
		Let $A\subset \Omega\setminus (I\times Y)$ be a Borel set, and let
		$$A_n=\{x\in A:|\mathring{w}''|(x)\ge \frac 1n\}.
		$$
		For $n$ large enough, $A_n$ is nonempty and $(A_n)\subset A$ monotonically converges to $A$. Then $0=
		(|\mathring{w}''|\mu)(A_n)\ge \frac1n\mu(A_n)\ge 0$, that implies  $\mu(A_n)=0$. In turn, we deduce $\mu(A)=0$.  Thus $\mu=\mu\mres (I\times Y)=\sum_{y\in Y} \mu\mres (I\times\{y\})$, since $Y$ is finite.
		For fixed $y\in Y$, let $\lambda_y$ be the Radon measure defined on the Borel sets of $I\times\{y\}$ by
		$$\lambda_y(B_y)=\mu\mres (I\times\{y\})(B_y)=\mu(B_y),
		$$
		for all Borel sets $B_y\subset I\times\{y\}$. For a Borel set $B\subset \Omega$, we have that
		$$\lambda_y\otimes\delta_y(B)=\lambda_y(B\cap \{x_2=y\})=\lambda_y(B\cap( I\times\{y\}))=\mu(B\cap( I\times\{y\}))=\mu\mres (I\times\{y\})(B)
		$$
		and hence $\mu\mres (I\times\{y\})=\lambda_y\otimes\delta_y.$
		%%%%\QED
	\end{proof}

\end{document}
